

\chapter{Jetson Nano MSR-Lab}



Der Jetson Nano des MSR-Labors befindet sich in einem Metallgehäuse, an dem ein Power- sowie Reset-Button eingelassen sind und die Kamera an der vorgesehenen Halterung befestigt ist (Abbildung~\ref{JetsonMSR1}).

\begin{figure}[H]
	\centering
	\includegraphics[height=0.4\textwidth, angle=-90]{Jetson/JetsonMSR1}
	\caption{Jetson Nano MSR-Lab im Case mit Pi-Kamera, Frontansicht} \label{JetsonMSR1}
\end{figure}

An der Seite des Jetson Nano (Abbildung~\ref{JetsonMSR2}) liegen die benötigten Anschlüsse - USB, HDMI, Ethernet, Micro USB sowie der microSD-Kartenslot - frei. 

\begin{figure}[H]
	\centering
	\includegraphics[height=0.4\textwidth, angle=-90]{Jetson/JetsonMSR2}
	\caption{Anschlüsse des Jetson Nano MSR-Lab im Case} \label{JetsonMSR2}
\end{figure}

An der Rückseite des Jetson Nano (Abbildung~\ref{JetsonMSR3}) befindet sich der J41-Header mit den GPIO-Pins.

\begin{figure}[H]
	\centering
	\includegraphics[height=0.4\textwidth, angle=-90]{Jetson/JetsonMSR3}
	\caption{J41-Header des Jetson Nano MSR-Lab im Case, Rückansicht} \label{JetsonMSR3}
\end{figure}